\documentclass[11pt]{article}

\usepackage[a4paper,margin=1in]{geometry}
\usepackage[T1]{fontenc}
\usepackage{lmodern}
\usepackage{microtype}
\usepackage{amsmath,amssymb,amsthm,mathtools}
\usepackage{booktabs,array}
\usepackage{xcolor}
\usepackage{enumitem}
\usepackage[hidelinks]{hyperref}

\newtheorem{theorem}{Theorem}[section]
\newtheorem{proposition}[theorem]{Proposition}
\newtheorem{lemma}[theorem]{Lemma}
\newtheorem{corollary}[theorem]{Corollary}
\theoremstyle{remark}
\newtheorem{remark}[theorem]{Remark}

\newcommand{\R}{\mathbb R}
\newcommand{\C}{\mathbb C}
\newcommand{\Sym}{\mathfrak S}
\newcommand{\one}{\mathbf 1}
\newcommand{\id}{I}
\newcommand{\HP}{\mathbb H}
\newcommand{\rank}{\operatorname{rank}}
\newcommand{\tr}{\operatorname{tr}}
\newcommand{\diag}{\operatorname{diag}}
\newcommand{\Span}{\operatorname{span}}
\newcommand{\partn}{\vdash}

\title{A Counterexample to Borcea--Br\"and\'en AIM Problem 36\\
\large A symmetric determinantal stable polynomial with a negative Schur coefficient}
\author{}
\date{Exact GPT-assisted certificate, July 2026}

\begin{document}
\maketitle

\begin{abstract}
Borcea--Br\"and\'en AIM Problem 36 asks whether every symmetric, homogeneous,
real-stable polynomial of degree $d$ in $n\ge d$ variables, with nonnegative
monomial coefficients, is Schur-positive.  We give an explicit counterexample
with $d=n=5$.  It is a determinant of a positive-definite linear matrix pencil,
so real stability and monomial nonnegativity are immediate.  Its exact Schur
expansion has coefficient $-2972$ at $s_{(1^5)}$.  We also give a conceptual
one-parameter construction from the Specht module $S^{(3,2)}$ and an exact
verification script using rational arithmetic only.  No floating-point
calculation enters the certificate.
\end{abstract}

\section{The conjecture and the counterexample}

A polynomial $f\in\R[x_1,\dots,x_n]$ is \emph{real stable} if
\[
  f(z_1,\dots,z_n)\ne0
  \qquad\text{whenever }z_1,\dots,z_n\in\HP
  :=\{z\in\C:\operatorname{Im}z>0\}.
\]
Problem 36 in the AIM workshop report asks whether a symmetric, homogeneous,
real-stable polynomial of degree $d\le n$, with nonnegative coefficients in the
ordinary monomial basis, must have nonnegative coefficients in the Schur basis
\cite[Problem~36]{AIM2007}.

Let $J_1,\dots,J_5$ be the following symmetric integer matrices:
\begingroup
\scriptsize
\[
J_1=
\begin{pmatrix}
38&19&19&2&36\\
19&38&2&19&36\\
19&2&38&19&36\\
2&19&19&38&36\\
36&36&36&36&72
\end{pmatrix},\qquad
J_2=
\begin{pmatrix}
38&19&19&17&21\\
19&38&17&19&21\\
19&17&38&34&36\\
17&19&34&38&36\\
21&21&36&36&42
\end{pmatrix},
\]
\[
J_3=
\begin{pmatrix}
38&34&19&17&36\\
34&38&17&19&36\\
19&17&38&19&21\\
17&19&19&38&21\\
36&36&21&21&42
\end{pmatrix},\qquad
J_4=
\begin{pmatrix}
38&4&19&2&21\\
4&8&2&4&6\\
19&2&38&4&21\\
2&4&4&8&6\\
21&6&21&6&42
\end{pmatrix},
\]
\[
J_5=
\begin{pmatrix}
8&4&4&2&6\\
4&38&2&19&21\\
4&2&8&4&6\\
2&19&4&38&21\\
6&21&6&21&42
\end{pmatrix}.
\]
\endgroup
Define
\begin{equation}
  q(x_1,\dots,x_5)
  :=\frac1{648}\det\!\left(\sum_{i=1}^5x_iJ_i\right).
  \label{eq:qdef}
\end{equation}

\begin{theorem}[Explicit counterexample]
\label{thm:explicit}
The polynomial $q$ in \eqref{eq:qdef} is symmetric, homogeneous of degree $5$,
real stable, and every one of its $126$ ordinary degree-$5$ monomial
coefficients is strictly positive.  Nevertheless,
\begin{align}
q={}&1728s_{(5)}+17712s_{(4,1)}+47115s_{(3,2)}
      +56643s_{(3,1,1)} \notag\\
   &\quad +62415s_{(2,2,1)}+56437s_{(2,1,1,1)}
      -2972s_{(1,1,1,1,1)}.
\label{eq:schurexpansion}
\end{align}
In particular, $q$ is not Schur-positive.
\end{theorem}

\begin{proof}
Homogeneity of degree $5$ is immediate from the determinant of a $5\times5$
linear pencil.  The successive leading principal minors of the matrices $J_i$
are
\[
\begin{array}{c|rrrrr}
 &\Delta_1&\Delta_2&\Delta_3&\Delta_4&\Delta_5\\ \hline
J_1&38&1083&28728&202176&1119744\\
J_2&38&1083&28728&202176&1119744\\
J_3&38&288&8208&202176&1119744\\
J_4&38&288&8208&46656&1119744\\
J_5&8&288&1728&46656&1119744.
\end{array}
\]
They are all positive, so Sylvester's criterion gives $J_i\succ0$ for every $i$.
If $z_i\in\HP$ and
$\left(\sum_i z_iJ_i\right)v=0$ for some nonzero $v\in\C^5$, then
\[
0=\operatorname{Im}\left\langle v,
       \left(\sum_i z_iJ_i\right)v\right\rangle
  =\sum_i \operatorname{Im}(z_i)\,\langle v,J_iv\rangle>0,
\]
a contradiction.  Thus $q$ is real stable.

An exact determinant expansion gives
\begin{align}
q={}&1728m_{(5)}+19440m_{(4,1)}+66555m_{(3,2)}
     +140910m_{(3,1,1)}\notag\\
 &\quad +250440m_{(2,2,1)}+547405m_{(2,1,1,1)}
     +1182860m_{(1,1,1,1,1)}.
\label{eq:monomialexpansion}
\end{align}
This proves both symmetry and strict positivity of every ordinary monomial
coefficient.

For completeness, order the partitions of $5$ as
\[
(5),(4,1),(3,2),(3,1,1),(2,2,1),(2,1,1,1),(1^5).
\]
With this order, the Kostka matrix in
$s_\lambda=\sum_\mu K_{\lambda\mu}m_\mu$ is
\[
K=
\begin{pmatrix}
1&1&1&1&1&1&1\\
0&1&1&2&2&3&4\\
0&0&1&1&2&3&5\\
0&0&0&1&1&3&6\\
0&0&0&0&1&2&5\\
0&0&0&0&0&1&4\\
0&0&0&0&0&0&1
\end{pmatrix}.
\]
If $c$ and $a$ are the row vectors of monomial-symmetric and Schur
coefficients, respectively, then $c=aK$.  Applying $K^{-1}$ to the coefficient
vector in \eqref{eq:monomialexpansion} gives \eqref{eq:schurexpansion}.  In
particular, the final coefficient can be checked from the single identity
\begin{align*}
[s_{(1^5)}]q
={}&c_{(5)}-2c_{(4,1)}-2c_{(3,2)}+3c_{(3,1,1)}
       +3c_{(2,2,1)}\\
&\qquad -4c_{(2,1,1,1)}+c_{(1^5)}\\
={}&1728-2(19440)-2(66555)+3(140910)+3(250440)\\
&\qquad -4(547405)+1182860=-2972.
\end{align*}
All polynomial expansions and the Kostka inversion are checked independently by
the exact script included with this note; see Section~\ref{sec:certificate}.
\end{proof}

\begin{corollary}
AIM Problems 36, 37, and 38 are false as stated.
\end{corollary}

\begin{proof}
Theorem~\ref{thm:explicit} disproves Problem 36.  For Problem 37, take
$\lambda=(1^5)$; since $f^{(1^5)}=1$, the proposed inequality would require
$-2972\ge1728$.  For Problem 38, take $\lambda=(5)$; since $f^{(5)}=1$, the
proposed inequality would require $1728\le-2972$.
\end{proof}

\section{Conceptual origin: a one-parameter Specht-module family}

The explicit matrices above are not a random numerical artifact.  They arise
from a simple equivariant construction on the Specht module $S^{(3,2)}$.

Let $M$ be the real permutation module with orthonormal basis
$\{e_{ab}:1\le a<b\le5\}$ indexed by the $2$-subsets of $[5]$.  Define an
$\Sym_5$-equivariant map
\[
B:M\longrightarrow\R^5,\qquad B(e_{ab})=e_a+e_b.
\]
Its matrix satisfies
\[
BB^{\mathsf T}=3I_5+\one\one^{\mathsf T},
\]
so $B$ is surjective and $V:=\ker B$ has dimension $5$.  Young's rule for the
permutation module on $2$-subsets gives
\[
M\cong S^{(5)}\oplus S^{(4,1)}\oplus S^{(3,2)},
\]
whereas the point-permutation module $\R^5$ is
$S^{(5)}\oplus S^{(4,1)}$.  Hence
\[
V\cong S^{(3,2)}.
\]
Let $\rho$ denote the orthogonal representation of $\Sym_5$ on $V$.

For $i\in[5]$, let $H_i\cong\Sym_4$ be the stabilizer of $i$, and set
\begin{equation}
P_i:=\frac12
\sum_{\substack{1\le a<b\le5\\a,b\ne i}}
\rho((ab)).
\label{eq:Pi}
\end{equation}

\begin{lemma}
Each $P_i$ is an orthogonal projection of rank $3$, and
\[
\rho(\sigma)P_i\rho(\sigma)^{-1}=P_{\sigma(i)}
\qquad(\sigma\in\Sym_5).
\]
Moreover, $\sum_{i=1}^5P_i=3I_V$.
\end{lemma}

\begin{proof}
The branching rule gives
\[
S^{(3,2)}\!\downarrow_{\Sym_4}
\cong S^{(3,1)}\oplus S^{(2,2)}.
\]
The sum of all six transpositions in $\Sym_4$ is central in $\R[\Sym_4]$.
On an irreducible Specht module $S^\mu$ it acts by the sum of the contents of
$\mu$.  These sums are
\[
\sum_{u\in(3,1)}c(u)=2,
\qquad
\sum_{u\in(2,2)}c(u)=0.
\]
Therefore the operator in \eqref{eq:Pi} acts as the identity on the
$3$-dimensional $S^{(3,1)}$ summand and as zero on the $2$-dimensional
$S^{(2,2)}$ summand.  It is thus an orthogonal rank-$3$ projection.  Covariance
follows by conjugating the transposition class sum.  Finally, $\sum_iP_i$
commutes with $\rho(\Sym_5)$, so Schur's lemma makes it scalar on the irreducible
space $V$.  Its trace is $5\cdot3=15$, hence it equals $3I_V$.
\end{proof}

For $t\ge0$, define
\[
A_i(t):=I_V+tP_i,
\qquad
p_t(x):=\det_V\!\left(\sum_{i=1}^5x_iA_i(t)\right).
\]
Each $A_i(t)$ is positive definite, with eigenvalues $1+t$ of multiplicity $3$
and $1$ of multiplicity $2$.

\begin{proposition}
For every $t\ge0$, the polynomial $p_t$ is symmetric, homogeneous of degree $5$,
real stable, and has nonnegative monomial coefficients.  Its Schur coefficients
are
\begin{align*}
[s_{(5)}]p_t&=(t+1)^3,\\
[s_{(4,1)}]p_t&=\frac{(t+1)^2(3t^2+8t+8)}2,\\
[s_{(3,2)}]p_t&=\frac{(t+1)(9t^4+48t^3+128t^2+160t+80)}{16},\\
[s_{(3,1,1)}]p_t&=\frac{(t+1)(9t^4+64t^3+168t^2+192t+96)}{16},\\
[s_{(2,2,1)}]p_t&=\frac{6t^5+41t^4+108t^3+156t^2+120t+40}{8},\\
[s_{(2,1,1,1)}]p_t&=\frac{6t^5+35t^4+96t^3+132t^2+96t+32}{8},\\
[s_{(1^5)}]p_t&=-\frac{9t^5-21t^4-56t^3-72t^2-48t-16}{16}.
\end{align*}
In particular,
\[
[s_{(1^5)}]p_5=-\frac{743}{2}<0.
\]
\end{proposition}

\begin{proof}
The covariance of the $A_i(t)$ proves symmetry.  Positive definiteness proves
real stability by the same imaginary-part argument used in
Theorem~\ref{thm:explicit}.  The coefficient of
$x_1^{\alpha_1}\cdots x_5^{\alpha_5}$ is a positive scalar multiple of the
mixed discriminant of the multiset containing $\alpha_i$ copies of $A_i(t)$.
Mixed discriminants of positive-semidefinite matrices are nonnegative; for
positive-definite arguments they are strictly positive.  Thus the monomial
coefficients are nonnegative (indeed positive).

The displayed Schur coefficients follow by exact determinant expansion and
Kostka inversion.  For an entirely transparent intermediate certificate, the
monomial-symmetric coefficients are listed in Appendix~\ref{app:family}.  The
included script derives the Kostka matrix by direct enumeration of semistandard
Young tableaux and independently checks every Schur coefficient by bialternant
extraction.
\end{proof}

\begin{remark}
The unique positive real zero of
$9t^5-21t^4-56t^3-72t^2-48t-16$ is approximately
$4.2907656153$.  Thus the construction gives an open ray of counterexamples,
not merely an isolated example.  This numerical value is not used anywhere in
the proof.
\end{remark}

\section{Relation between the conceptual and explicit certificates}

The script chooses a rational basis $K$ of $V=\ker B$ and lets
$G=K^{\mathsf T}K$ be its Gram matrix.  In that basis, the matrices of $P_i$ are
$G$-self-adjoint.  At $t=5$, set
\[
J_i:=2G(I+5P_i).
\]
Then $J_i$ is an ordinary symmetric positive-definite integer matrix and
\[
\det\!\left(\sum_i x_iJ_i\right)
=2^5\det(G)\,p_5(x)=5184p_5(x).
\]
Consequently,
\[
q(x)=\frac1{648}\det\!\left(\sum_i x_iJ_i\right)=8p_5(x),
\]
which explains the normalization in \eqref{eq:qdef}.  The matrices displayed in
Section~1 are exactly the resulting $J_i$.

\section{Exact computational certificate}
\label{sec:certificate}

The accompanying file \path{verify_problem36_counterexample.py} requires only Python and SymPy.
It uses exact rational and integer arithmetic throughout.  Running
\begin{verbatim}
python verify_problem36_counterexample.py --verbose
\end{verbatim}
performs all of the following checks:
\begin{enumerate}[leftmargin=2em]
\item constructs $V=\ker B$ and the five projections $P_i$;
\item verifies $P_i^2=P_i$, $G$-self-adjointness, rank $3$, full
      $\Sym_5$-covariance, and $\sum_iP_i=3I$;
\item constructs the five matrices $J_i$ and verifies positive definiteness by
      exact principal minors;
\item expands the determinant into all $126$ degree-$5$ monomials and verifies
      constancy on every $\Sym_5$-orbit and strict positivity;
\item computes the Kostka matrix by enumerating semistandard Young tableaux,
      rather than importing or hard-coding it;
\item obtains the Schur expansion by exact matrix inversion;
\item independently recomputes every Schur coefficient using the bialternant
      extraction formula;
\item verifies the full one-parameter coefficient formulas stated above.
\end{enumerate}
The optional command
\begin{verbatim}
python verify_problem36_counterexample.py \
       --write-json counterexample_certificate.json
\end{verbatim}
produces the machine-readable certificate included in the package.

\appendix
\section{Monomial coefficients of the one-parameter family}
\label{app:family}

Writing $p_t=\sum_{\mu\partn5}c_\mu(t)m_\mu$, exact expansion gives
\begin{align*}
c_{(5)}(t)&=(t+1)^3,\\
c_{(4,1)}(t)&=\frac{(t+1)^2(3t^2+10t+10)}2,\\
c_{(3,2)}(t)&=\frac{(t+1)(9t^4+72t^3+232t^2+320t+160)}{16},\\
c_{(3,1,1)}(t)&=\frac{(t+1)(t+2)(9t^3+62t^2+120t+80)}8,\\
c_{(2,2,1)}(t)&=\frac{39t^5+317t^4+1040t^3+1728t^2+1440t+480}{16},\\
c_{(2,1,1,1)}(t)&=\frac{45t^5+348t^4+1100t^3+1764t^2+1440t+480}{8},\\
c_{(1^5)}(t)&=\frac{99t^5+765t^4+2320t^3+3600t^2+2880t+960}{8}.
\end{align*}
All coefficients in these numerators are positive.  Combining these formulas
with the last column of $K^{-1}$,
\[
(1,-2,-2,3,3,-4,1)^{\mathsf T},
\]
gives
\[
[s_{(1^5)}]p_t
=-\frac{9t^5-21t^4-56t^3-72t^2-48t-16}{16}.
\]

\begin{thebibliography}{9}
\bibitem{AIM2007}
J.~Borcea, P.~Br\"and\'en, G.~Csordas, and V.~Vinnikov (organizers),
\emph{P\'olya--Schur--Lax problems: hyperbolicity and stability preservers},
AIM Workshop Summary, American Institute of Mathematics, 2007,
\url{https://www.aimath.org/pastworkshops/polyaschurlaxrep.pdf}.
\end{thebibliography}

\end{document}
